\subsection{R 上的矩阵空间的拓扑}

域 \(\mathbf{R}\) 上向量空间的一个重要例子是空间 \(\mathbf{M}_{m,n}(\mathbf{R})\) ——由 \(m\) 行 \(n\) 列、元素属于 \(\mathbf{R}\) 的矩阵组成的空间；这是 \(mn\) 维的空间（在 \(\mathbf{R}\) 上），因此同胚于 \(\mathbf{R}^{mn}\)。由 § 5 的命题 5，乘积 \(X\) . \(\Upsilon\) （两个矩阵 \(X \in \mathbf{M}_{m,n}(\mathbf{R})\) 与 \(\Upsilon \in \mathbf{M}_{n,p}(\mathbf{R})\) 的乘积）是 \((X, \Upsilon)\) 的连续函数。特别地，方阵空间 \(\mathbf{M}_n(\mathbf{R})\) （\(n\) 阶的）的拓扑与 \(\mathbf{M}_n(\mathbf{R})\) 上的环结构相容。此外：

命题 6. 在环 \(\mathbf{M}_n(\mathbf{R})\) 中，非奇异矩阵组成的群 \(\mathbf{GL}_n(\mathbf{R})\) 是稠密开子集，而且在这个集上诱导的拓扑与其群结构相容。

（*）若 E, F, G 是域 K 上三个向量空间，则称 \(E \times F\) 到 G 中的映射 f 为双线性的，如果恒有

\[
\begin{array}{r l} f (x + x ^ {\prime}, y) & = f (x, y) + f (x ^ {\prime}, y), f (x, y + y ^ {\prime}) = f (x, y) + f (x, y ^ {\prime}), \\ f (\lambda x, y) & = f (x, \lambda y) = \lambda f (x, y) \end{array}
\]

若 X 是非奇异方阵，则 \(X^{-1}\) 的元素是 X 的元素的有理函数；因此这些函数在 X 的一个邻域中有定义且连续，从而这个邻域中的每个矩阵 Y 都是非奇异的，而且映射 \(Y \to Y^{-1}\) 在点 X 连续；因此 \(\mathbf{GL}_{n}(\mathbf{R})\) 在 \(\mathbf{M}_{n}(\mathbf{R})\) 中是开的，并且 \(\mathbf{GL}_{n}(\mathbf{R})\) 的拓扑与其群结构相容。

最后，\(\mathbf{GL}_{n}(\mathbf{R})\) 是行列式为零的方阵 X 所成的集的补集；由于 X 的行列式是 X 的元素的多项式，第 4 小节的命题 3 表明 \(\mathbf{GL}_{n}(\mathbf{R})\) 在 \(\mathbf{M}_{n}(\mathbf{R})\) 中稠密。

